% Generated by tools/textbook/export_listings.js from tools/concepts/assessments.json: do not edit.
\begin{chapstatus}{The thread of this chapter}
A register is flip-flops that sample on a clock edge: this is how the counter of Lecture 02 survived from one frame to the next, in hardware. Next: how registers exchange data.
\par\medskip\textbf{Exit ticket} (two to four minutes, alone, on paper):
\begin{enumerate}[nosep,leftmargin=2em]
\item The kit's clock runs at 160 MHz. What is its period, and why must every combinational network between two registers settle within it?
\item What is the difference between a D latch and a D flip-flop, and which one lets a register be read and written in the same clock cycle safely?
\end{enumerate}
\end{chapstatus}
